% 15-minute presentation. Compile from this folder: pdflatex slides.tex (twice).
% Plain Beamer, 16:9, large type, one figure per slide. Figures come from ../experiments/.
\documentclass[aspectratio=169,14pt]{beamer}
\usepackage[T1]{fontenc}
\usepackage{lmodern}
\usepackage{booktabs}
\usepackage{amsmath}
\graphicspath{{../experiments/}}

% ---- minimal style -------------------------------------------------
\definecolor{ink}{HTML}{1F2933}
\definecolor{accent}{HTML}{B7410E}   % burnt orange
\definecolor{muted}{HTML}{6B7280}
\setbeamertemplate{navigation symbols}{}
\setbeamercolor{normal text}{fg=ink,bg=white}
\setbeamercolor{frametitle}{fg=ink}
\setbeamercolor{structure}{fg=accent}
\setbeamercolor{alerted text}{fg=accent}
\setbeamerfont{frametitle}{series=\bfseries,size=\Large}
\setbeamertemplate{frametitle}{\vspace*{0.4em}\insertframetitle\par\vspace{-0.1em}%
  {\color{accent}\rule{2.2em}{2pt}}\par}
\setbeamertemplate{itemize items}{\color{accent}\textbullet}
\setbeamertemplate{footline}{\hfill{\color{muted}\scriptsize\insertframenumber}\hspace{1.2em}\vspace{0.8em}}
\setlength{\leftmargini}{1.1em}

\newcommand{\bignum}[2]{{\color{accent}\bfseries\huge #1}\\[-0.1em]{\color{muted}\small #2}}
\newcommand{\sub}[1]{{\color{muted}\small #1}}
\newcommand{\fig}[2][0.74]{\centering\includegraphics[height=#1\textheight,width=\textwidth,keepaspectratio]{#2}}

\begin{document}

% ---- 1 -------------------------------------------------------------
\begin{frame}[plain]
  \vspace{1.5em}
  {\bfseries\LARGE Lightweight Temperature Predictor}\\[0.3em]
  {\Large\color{accent} Calibrate once, predict fast}\\[1.6em]
  Ethan Bobby Kurien \quad Kaivalya Mase\\[0.2em]
  \sub{Computer Architecture course project \textperiodcentered{} Thermal management}
\end{frame}

% ---- 2 -------------------------------------------------------------
\begin{frame}{The problem}
  \begin{itemize}
    \item Hot cores slow down, age and fail.
    \item DTM throttles them, and needs temperatures \alert{every millisecond}, ideally a few ms
          ahead.
    \item HotSpot is accurate but slow:
  \end{itemize}
  \vspace{0.6em}
  \begin{columns}[T]
    \column{0.45\textwidth}\centering \bignum{0.24 s}{per steady answer}
    \column{0.45\textwidth}\centering \bignum{5.5 ms}{per 1 ms transient step}
  \end{columns}
\end{frame}

% ---- 3 -------------------------------------------------------------
\begin{frame}{Our approach: two stages}
  \begin{columns}[T]
    \column{0.47\textwidth}
      {\bfseries\color{accent} 1. Calibrate} \sub{(once per floorplan)}\\[0.4em]
      Run HotSpot with 1\,W in one core at a time.\\[0.3em]
      8 runs for steady, 24 for transient.\\[0.3em]
      Store the coefficients: \alert{8\,kB}.
    \column{0.47\textwidth}
      {\bfseries\color{accent} 2. Predict} \sub{(many times)}\\[0.4em]
      Power of each core in, temperatures out.\\[0.3em]
      No HotSpot, a few multiplications.\\[0.3em]
      \alert{About 3\,$\mu$s} per answer.
  \end{columns}
\end{frame}

% ---- 4 -------------------------------------------------------------
\begin{frame}{Key insight: HotSpot is linear}
  Double the power, every rise doubles. Rises from different cores add up.\\
  \sub{Checked to $10^{-6}$\,K, steady and transient (P0-2).}
  \vspace{0.8em}
  \[
    T = T_\text{amb} + G\,P
  \]
  \vspace{-0.2em}
  \begin{itemize}
    \item $G_{ij}$: how much core $j$'s watt heats block $i$.
    \item 1\,W in core $j$ alone $\Rightarrow$ column $j$ of $G$.
    \item So the steady model is \alert{exact}, not approximate.
  \end{itemize}
\end{frame}

% ---- 5 -------------------------------------------------------------
\begin{frame}{Setup}
  \begin{columns}[c]
    \column{0.40\textwidth}
      \begin{itemize}
        \item 23 TA floorplans: 2 big + 6 little cores
        \item HotSpot $64\times64$ grid
        \item We patched HotSpot:
          \begin{itemize}
            \item 6 decimals, not 0.01\,K
            \item \alert{29$\times$} faster transients
            \item memory leak fixed
          \end{itemize}
      \end{itemize}
    \column{0.58\textwidth}
      \fig[0.86]{P0-5d_generated_floorplans/results/examples.png}
  \end{columns}
\end{frame}

% ---- 6 -------------------------------------------------------------
\begin{frame}{Calibration: heat moves at many speeds}
  \fig[0.58]{P1-1_step_calibration/results/step_response_fits.png}\\[0.2em]
  \sub{Time constants from 1\,ms (silicon) to 8\,s (heat sink). One exponential fails (14\,K);
  \alert{four} fit within 0.46\,K.}
\end{frame}

% ---- 7 -------------------------------------------------------------
\begin{frame}{Is calibration really needed?}
  \begin{columns}[c]
    \column{0.55\textwidth}
      \fig[0.78]{P1-3_is_calibration_needed/results/error_by_method.png}
    \column{0.43\textwidth}
      Same cores, same power, only the layout differs.\\[0.6em]
      Without calibration:\\
      \alert{0.5--2.7\,K} mean, up to \alert{12\,K}.\\[0.6em]
      With 8 runs (4\,s):\\
      \alert{0.00003\,K}.
  \end{columns}
\end{frame}

% ---- 8 -------------------------------------------------------------
\begin{frame}{Steady predictor (S1)}
  \vspace{0.6em}
  \begin{columns}[T]
    \column{0.40\textwidth}\centering \bignum{\mbox{0.00003\,K}}{worst error}
    \column{0.28\textwidth}\centering \bignum{\mbox{2.6\,$\mu$s}}{per prediction}
    \column{0.28\textwidth}\centering \bignum{\mbox{\textless\,1\,kB}}{model}
  \end{columns}
  \vspace{1.4em}
  \begin{itemize}
    \item 90{,}000$\times$ faster than HotSpot.
    \item Always picks the right hottest core.
    \item Tested on 7{,}248 unseen HotSpot runs.
  \end{itemize}
\end{frame}

% ---- 9 -------------------------------------------------------------
\begin{frame}{The hot spot inside a core}
  \begin{columns}[c]
    \column{0.55\textwidth}
      \fig[0.78]{P2-3_hot_spots/results/hot_spot_error.png}
    \column{0.43\textwidth}
      Core average misses the peak by up to \alert{9\,K}.\\[0.6em]
      The hot spot moves with the workload.\\[0.6em]
      \alert{S1-hot-max}: track a few candidate cells, take the max.\\
      0.08\,K mean, 3\,kB, no extra runs.
  \end{columns}
\end{frame}

% ---- 10 ------------------------------------------------------------
\begin{frame}{Transient predictor (T3)}
  \begin{columns}[c]
    \column{0.58\textwidth}
      \fig[0.78]{P3-1_transient_predictor/results/trace_example.png}
    \column{0.40\textwidth}
      Sum of 4 exponentials per pair of cores.\\[0.6em]
      \alert{0.06--0.09\,K} mean\\
      0.3\,K worst\\[0.6em]
      3\,$\mu$s per step, 4\,kB\\
      1{,}800$\times$ faster than HotSpot
  \end{columns}
\end{frame}

% ---- 11 ------------------------------------------------------------
\begin{frame}{Forecasting ahead for DTM}
  \begin{columns}[c]
    \column{0.55\textwidth}
      \fig[0.78]{P3-2_forecast_horizon/results/error_vs_horizon.png}
    \column{0.43\textwidth}
      ``Stays the same'' (no model): up to \alert{8\,K} wrong 10\,ms ahead.\\[0.6em]
      T3 anchored on the sensors: \alert{0.004\,K} mean 10\,ms ahead.
  \end{columns}
\end{frame}

% ---- 12 ------------------------------------------------------------
\begin{frame}{Accuracy vs size}
  \begin{columns}[c]
    \column{0.64\textwidth}
      \fig[0.82]{P3-3_pareto/results/pareto.png}
    \column{0.34\textwidth}\small
      T3 is the best accuracy per kilobyte from 2 to 6\,kB.\\[0.6em]
      Each extra exponential: +1\,kB, error $\div$ 2--4.\\[0.6em]
      Cross-check: HotSpot's own 16{,}396-node network gives the same $G$.
  \end{columns}
\end{frame}

% ---- 13 ------------------------------------------------------------
\begin{frame}{Does it hold up?}
  \begin{columns}[c]
    \column{0.52\textwidth}
      \fig[0.74]{extended_validation/results/error_by_style.png}
    \column{0.46\textwidth}\small
      \begin{itemize}
        \item 52 new floorplans: same accuracy.
        \item 3$\times$ power, bursty, migration: fine.
        \item New ambient: free.
        \item New heat sink: recalibrate.
        \item Leakage: 0.1\,K, fixed by a short loop.
        \item Other time step: 2\,K. Calibrate at the DTM step.
      \end{itemize}
  \end{columns}
\end{frame}

% ---- 14 ------------------------------------------------------------
\begin{frame}{Proactive DTM}
  \begin{columns}[c]
    \column{0.58\textwidth}
      \fig[0.78]{P6-2_dtm/results/dtm_s1_fp001_tight.png}
    \column{0.40\textwidth}
      HotSpot as the chip, controller every 10\,ms.\\[0.6em]
      \alert{60--130$\times$} less time over the limit than reactive DTM.\\[0.6em]
      Cost: 1--1.6\,\% throughput.
  \end{columns}
\end{frame}

% ---- 15 ------------------------------------------------------------
\begin{frame}{Scorecard}
  \centering\small
  \begin{tabular}{@{}l l c c@{}}
    \toprule
    & Metric & Target & Achieved \\
    \midrule
    S1 & worst error & $\le$ 0.2\,K & \alert{0.00003\,K} \\
       & time & $\le$ 20\,$\mu$s & \alert{2.6\,$\mu$s} \\
       & size & $\le$ 5\,kB & \alert{464--928\,B} \\
    \midrule
    T3 & mean error & $\le$ 0.5\,K & \alert{0.06--0.09\,K} \\
       & time per step & $\le$ 50\,$\mu$s & \alert{3.1\,$\mu$s} \\
       & size & $\le$ 50\,kB & \alert{4\,kB} \\
       & calibration time & $\le$ 2\,min & 2.75\,min \\
    \bottomrule
  \end{tabular}
\end{frame}

% ---- 16 ------------------------------------------------------------
\begin{frame}{Takeaways}
  \begin{itemize}
    \item HotSpot is linear, so a tiny model can be exact.
    \item 4 seconds of calibration beats every shortcut.
    \item Layout matters: edges, neighbours, die size.
    \item Averages hide hot spots; track the peak.
    \item With a forecast, DTM avoids overheating instead of fixing it.
  \end{itemize}
\end{frame}

% ---- 17 ------------------------------------------------------------
\begin{frame}[plain]
  \centering\vfill
  {\bfseries\LARGE Thank you}\\[0.6em]
  {\color{accent}\Large Questions?}
  \vfill
\end{frame}

% ---- backup --------------------------------------------------------
\appendix
\begin{frame}[fragile]{Backup: the tool}
  \footnotesize
\begin{verbatim}
python3 -m src.ltp calibrate s4_fp001 -o model.npz
python3 -m src.ltp steady    model.npz --power 20,20,8,8,8,8,8,8
python3 -m src.ltp transient model.npz --ptrace trace.ptrace -o out.ttrace
python3 -m src.ltp bench     model.npz s4_fp001
\end{verbatim}
  \normalsize
  \sub{bench on s4\_fp001: 8{,}032\,B model; steady $5\times10^{-6}$\,K at 2.6\,$\mu$s;
  transient 0.051\,K mean at 3.5\,$\mu$s per step.}
\end{frame}

\begin{frame}{Backup: what S1 sees vs HotSpot}
  \fig[0.72]{extended_validation/results/heatmaps.png}
\end{frame}

\end{document}
